\chapter{Introduction}
\label{cap:introduction}


Over the past few decades, mobile applications have experienced ex-
plosive growth, increasingly replacing traditional desktop software
and web browsing. This transformation is not only technological but
also cultural, reflecting the shift toward a digital society [22]—one
in which everyday activities, social interactions, and information
flows are deeply embedded in digital infrastructure and mobile tech-
nologies. This transformation is driven by several key factors: the
widespread adoption of mobile devices [15, 31], their growing com-
putational power [13], increased battery life [8], and an expansion

of use cases [19]. Modern smartphones offer a wider range of func-
tionalities compared to conventional computers, largely due to their
built-in sensors—such as gyroscopes, GPS modules, and light sen-
sors—and actuators like vibration motors and cameras [33]. These
features have enabled a richer, more interactive user experience
and opened the door to innovative applications in diverse domains,
including health [4], education [21], and serious games [23]. Ap-
plications that interact with the physical world—through sensors
or cameras—and mobile games have become particularly preva-
lent [27]. Prominent examples such as Google Maps, Waze, Strava,
Snapchat, and Instagram illustrate the popularity and ubiquity of
such mobile apps, each boasting millions of downloads.
The shift from browser-based to mobile-based internet use has
been striking. In 2011, only 6% of users accessed the internet via
mobile phones; today, that figure exceeds 93%, representing ap-
proximately 4 billion people [11]. This trend is mirrored in social
media usage: over 80% of Facebook users now access the platform
exclusively through mobile devices, a dramatic increase from 30%
in 2014 [20, 30]. As a result, many companies are prioritizing mo-
bile platforms to deliver personalized, engaging experiences. To
achieve this, they increasingly rely on sensor data to infer user be-
havior and preferences. For instance, GPS data can reveal movement
patterns [29], while gyroscopes can detect specific gestures or ac-
tions [7]. However, this data collection raises several challenges, in-
cluding faster battery depletion [14], higher data consumption [25],
and significant privacy concerns [10].
As mobile hardware continues to evolve, there is an increasing
convergence between physical and digital spaces. Location-aware
technologies now support a wide spectrum of applications, from
real-time navigation and location-based recommendations to aug-
mented reality games and immersive media experiences. This shift
has transformed how individuals interact not only with their de-
vices but also with their surroundings, turning smartphones into
context-sensitive tools capable of enhancing real-world interac-
tions.
Despite this evolution, many existing location-based applica-
tions remain proprietary, rigid in functionality, or focused pre-
dominantly on data consumption rather than creation. There is a
notable lack of open, extensible platforms that empower users to
actively contribute multimedia content—such as audio, images, or
video—anchored to specific physical locations. Furthermore, tools
that enable users to collaboratively create and explore geo-tagged
content in a transparent and privacy-conscious environment are
scarce. These gaps present a clear opportunity for developing plat-
forms that are both technically flexible and socially engaging.
Within this context, we introduce GeoMedia: an open-source mo-
bile platform that enables users to share comments and multimedia
content—such as photos, audio, or video—directly on a geographicGoodIT ’25, September 03–05, 2025, Antwerp, Belgium
Alberto Morini, Lorenzo Perinello, and Claudio E. Palazzi
map, linking each post to the user’s current location. GeoMedia
aims to support interactive spatial storytelling, collaborative ex-
ploration, and personalized geolocation-based content sharing. Its
modular architecture allows developers and researchers to extend
the platform for various use cases, from tourism and urban art to
education and field research.
The remainder of this paper is organized as follows: Section 2 re-
views relevant related work in the domain of location-based social
networks and geolocation in mobile applications. Section 3 intro-
duces the design and main features of the GeoMedia platform. Sec-
tion 4 describes the system architecture, detailing the presentation,
application, and data layers. Section 5 focuses on the development
of the mobile application, including its main components, such as
the map interface, exclusive posts, and media handling. Section 6
presents the backend server implementation and its integration
with the database, while Section 7 outlines the database schema
and the use of stored procedures to manage data efficiently. Finally,
Sections 8 and 9 present directions for future work and conclude


% \noindent Esempio di utilizzo di un termine nel glossario \\
% \gls{api}. \\

% \noindent Esempio di citazione in linea \\
% \cite{site:agile-manifesto}. \\

% \noindent Esempio di citazione nel pie' di pagina \\
% citazione\footcite{womak:lean-thinking} \\

% \section{L'azienda}

% Descrizione dell'azienda.

% \section{L'idea}

% Introduzione all'idea dello stage.

% \section{Organizzazione del testo}

% \begin{description}
%     \item[{\hyperref[cap:processi-metodologie]{Il secondo capitolo}}] descrive ...
    
%     \item[{\hyperref[cap:descrizione-stage]{Il terzo capitolo}}] approfondisce ...
    
%     \item[{\hyperref[cap:analisi-requisiti]{Il quarto capitolo}}] approfondisce ...
    
%     \item[{\hyperref[cap:progettazione-codifica]{Il quinto capitolo}}] approfondisce ...
    
%     \item[{\hyperref[cap:verifica-validazione]{Il sesto capitolo}}] approfondisce ...
    
%     \item[{\hyperref[cap:conclusioni]{Nel settimo capitolo}}] descrive ...
% \end{description}

% Riguardo la stesura del testo, relativamente al documento sono state adottate le seguenti convenzioni tipografiche:
% \begin{itemize}
% 	\item gli acronimi, le abbreviazioni e i termini ambigui o di uso non comune menzionati vengono definiti nel glossario, situato alla fine del presente documento;
% 	\item per la prima occorrenza dei termini riportati nel glossario viene utilizzata la seguente nomenclatura: \emph{parola}\glsfirstoccur;
% 	\item i termini in lingua straniera o facenti parti del gergo tecnico sono evidenziati con il carattere \emph{corsivo}.
% \end{itemize}
